\section{Introduction}\label{sec:intro}

Soft sensors estimate quality variables that are measured slowly, and plants correct them against
laboratory or analyzer values that arrive late \citep{fortuna2007,kadlec2009}: the debutanizer's
gas chromatograph reports about 45~min after sampling \citep[Sec.~6.2]{fortuna2007}, and the
Tennessee Eastman product analyzer after 0.25~h \citep{downs1993}. Soft-sensor research treats
such labels as delayed, infrequent or irregular measurements \citep{guo2014,shardt2016}. An
operator also needs to know how far to trust the estimate, and adaptive conformal inference
\citep{gibbs2021aci} is a natural way to say it: an interval around the prediction whose level is
adjusted online so that a target share of labels falls outside. On a plant, those adjustments can
only use the late labels.

The adjustment is a feedback loop, and the delay puts dead time into it. With the step size tuned
on history that ignores the delay, the loop on the Tennessee Eastman process issued 6.9\,\% of its
intervals uninformative (infinite, or wider than the range of the measured variable) at the
analyzer's documented delay, 6.0 percentage points more than without delay, while the coverage
that is usually reported read 0.897 (Section~\ref{sec:results:diagnosis}). An operator cannot act
on such an interval, and the coverage figure does not show it.

Online conformal methods with delayed feedback exist. Adaptive conformal inference has been run on
labels that arrive several steps late \citep{dixit2023,szabadvary2024}, and at the common step of
0.005 it issued unbounded intervals in four of five electricity-market regions
\citep{wang2026cf}; its coverage under delay
has been bounded, with the step chosen per delay by experiment \citep{elhalabi2026}; interleaved
chains of delay-free updates avoid the dead time \citep{elhalabi2026,cherenson2026,li2026zhu}; and
AcMCP (autocorrelated multi-step conformal prediction) adds a forecast of the multi-step error to
a delayed loop \citep{wang2024acmcp}, as conformal PID control forecasts the score quantile
\citep{angelopoulos2023pid} and SPCI (sequential predictive conformal inference) the residual
quantiles \citep{xu2023spci}. The step size has been studied without delay \citep{zaffran2022,
angelopoulos2024decay}, and betting removes it altogether \citep{podkopaev2024}; neither that
paper nor this one examines betting under delay. Online learning with delayed feedback has a
regret theory of its own \citep{joulani2013,quanrud2015}. On process and prognostics data, a
PID-guided conformal
soft sensor uses each label at once \citep{meng2026}, and delayed-feedback conformal prediction
for fuel-cell prognostics raises its step with a label-free uncertainty signal \citep{qiu2026}.
What a plant lacks is an analysis of the conformal loop under laboratory delay: the limit that the
delay sets on the step, why the persistence of process residuals sets the usable step far below
that limit, how to tune the step so that it holds a budget out of sample, and what width a better
analyzer would buy.

This paper supplies that analysis and tests it. (i) The calibrator is an integral loop with dead
time: its linearisation is stable only for steps below $2\sin(\pi/(4\theta+2))$, and its marginal
coverage converges to nominal however many intervals are uninformative
(Section~\ref{sec:analysis}). (ii) Replaying the calibration residuals through the delayed loop
selects a step that holds a budget for the uninformative share (Section~\ref{sec:remedies:replay}).
(iii) Predicting the residual over the dead time, as a minimum-variance controller predicts a
disturbance \citep{astrom1970}, narrows the interval at the same budget, makes the misses less
persistent and admits a larger step (Section~\ref{sec:remedies:comp}). (iv) The width that this
compensation buys follows
from the autocorrelation of the calibration residuals before deployment, so the choice of analyzer
can be read off a chart (Section~\ref{sec:remedies:p3}). Each claim was tested against a rule
fixed with its threshold before any result of the study existed, in a pre-registered study on 15
Tennessee Eastman runs and the debutanizer data, against four baseline configurations
(Sections~\ref{sec:study} and~\ref{sec:results}).

The paper does not claim the conformal soft sensor, nor forecasting at the label horizon, both of
which are prior art; it claims their analysis and remedy under laboratory delay. The diagnosis was
shown on the simulated plant and not on the debutanizer; where the evaluation rows resembled the
calibration rows, compensated split conformal matched the delayed loop; and the predicted width
ratio is conservative on average, not a guarantee. An earlier preprint ran a delayed loop on both data sets
with the step size and the delay set by convention \citep{busico_m1}; here both follow from the
plant. Section~\ref{sec:loop} sets out the loop and the documented delays,
Section~\ref{sec:analysis} the analysis and the remedies, Section~\ref{sec:study} the study and
Section~\ref{sec:results} its results; Section~\ref{sec:discussion} turns them into guidance, and
Section~\ref{sec:conclusion} concludes.
